% Definicion del problema
\section{Problem definition}
\label{sec:problem}

We follow the problem definition of \citet{Paquay2016} with the extensions of
\citet{Hessler2025}. The length, width and height of a ULD correspond to the $x$, $y$ and $z$
axes, respectively; $z$ is the vertical axis.

\paragraph{Boxes.}
A set $I$ of rectangular boxes must be loaded. Each box $i \in I$ has a length $l_i$, a width
$w_i$, a height $h_i$ and a weight $q_i$. The boxes are placed orthogonally, and each box has a
set of allowed orientations: for each of its three dimensions, the instance states whether the
box can be placed with that dimension in the vertical direction (each allowed vertical
dimension gives two orientations, rotated or not around the vertical axis). A box can be
fragile, in which case nothing can be placed on top of it. The weight of a box is uniformly
distributed, so that its CoG is its geometric center. The instance files of \citet{Paquay2018b}
also flag very dense boxes; as \citet{Hessler2025}, we do not use this flag.

\paragraph{ULDs.}
A set $T$ of ULD types is available, in unlimited number or with a given number $n_k$ of ULDs
of each type $k \in T$. A ULD of type $k$ is the bounding box $[0, L_k] \times [0, W_k] \times
[0, H_k]$ from which up to four corners are cut by planes orthogonal to the $x$-$z$ plane, to fit
the fuselage of the aircraft. We represent each cut by a half-plane $\{(x, z) :
a\,x + b\,z \le c\}$ in the $x$-$z$ plane that contains the usable space; a box is inside the ULD
if it is inside its bounding box and its projection on the $x$-$z$ plane satisfies all the cuts.
The ULD type also defines its usable volume $V_k$, its weight capacity $Q_k$, the maximum
deviations $d^x_k$ and $d^y_k$ of the CoG of the load from the geometric center in the length and
width directions, and the maximum height $d^z_k$ of the CoG. Following \citet{Hessler2025}, the
floor of the ULD may have an edge of width $e$ for the cargo net locks: no box can be loaded
over the edge below the height $\delta$ (the vertical edge offset). To load large boxes at the
bottom, a substructure of height $\delta$ can be used, which raises the floor (without the edge)
to the height $\delta$; in that case no box can be loaded below the height $\delta$.

\paragraph{Constraints.}
A loading plan assigns each loaded box to a ULD, an orientation and a position. It is feasible
if:
\begin{enumerate}
\item each box is inside its ULD and, if the ULD has an edge, it does not overlap the edge
	below the height $\delta$;
\item the boxes do not overlap;
\item the total weight of the boxes in each ULD of type $k$ does not exceed $Q_k$;
\item each box whose base is not on the floor is supported: either the four corners of its base
	rest on other boxes (or on the floor or the substructure), or part of its base rests directly
	on other boxes and at least a proportion $o$ of its base area rests directly on boxes or on
	padding material \citep{Hessler2025}; padding material can fill gaps of height at most
	$\hbar$ between the base of a box and the top of the boxes below;
\item no box, and no padding material, is placed on top of a fragile box;
\item the CoG $(\bar x, \bar y, \bar z)$ of the load of each ULD satisfies
	$|\bar x - L_k/2| \le d^x_k$, $|\bar y - W_k/2| \le d^y_k$ and $\bar z \le d^z_k$.
\end{enumerate}
The original problem of \citet{Paquay2016} corresponds to $e = \delta = \hbar = 0$ and $o = 1$,
that is, a box is supported only if the four corners of its base are supported.
\citet{Hessler2025} do not consider the vertical CoG constraint; we consider it in all the
experiments. Moreover, both \citet{Paquay2018b} and \citet{Hessler2025} treat the horizontal
CoG constraint as a soft constraint, since their solutions contain ULDs with the CoG out of
range; we report the number of such ULDs for all the methods.

\paragraph{Objective.}
When the number of ULDs is unlimited, all the boxes must be loaded and the objective is to
minimize the number of ULDs. When it is limited, some boxes may be left unloaded, and the
objective is first to maximize the number of loaded boxes and then to minimize the number of
ULDs. These are the criteria used in the instance-by-instance comparison of
Section~\ref{sec:experiments}, and the ones reported by \citet{Hessler2025}.
\citet{Paquay2016} formulate the objective as the minimization of the unused volume, that is,
of the total volume of the used ULDs, and \citet{Hessler2025} maximize the loaded volume and then
minimize the total volume of the ULDs. Our algorithm minimizes the number of ULDs, and at the end
it tries to move the boxes of its least utilized ULD into a smaller one
(Section~\ref{sec:alg-multi}).
\todo{reportar el volumen total de los ULDs de nuestras soluciones y del RGS.}
